% Zumdahl Chapter 6 -- Thermochemistry. ELABORATED notes, 4 block lessons.
\documentclass{apchem-notes}

\apcunitword{Chapter}
\apcunit{6}
\apcunittitle{Thermochemistry}
\apcdoctitle{Guided Notes}
\apcsubtitle{Zumdahl \S6.1--6.4 \textbullet\ PDF pp.~282--308 \textbullet\ 4 blocks}

\begin{document}
\apctitleblock

\begin{apcnote}[How this chapter fits the AP course]
This is the \textbf{cleanest chapter-to-unit match in the book}. Sections
6.1--6.4 are essentially all of CED Unit 6, in the textbook's own order ---
you can read straight through without untangling anything.

Unit 6 adds only two topics from elsewhere, and you have already met both:
\textbf{energy of phase changes} (CED 6.5, Zumdahl \S10.8) came up in Unit 3
Block 2, and \textbf{bond enthalpies} (CED 6.7, Zumdahl \S8.8) in Unit 2
Block 1. Treat those as callbacks to your own notes.

\S6.5--6.6 (present and new energy sources) are skipped --- off-syllabus.
\end{apcnote}

% =====================================================================
\blocklesson{The Nature of Energy}{Zumdahl \S6.1}

\begin{objectives}
  \item Define system and surroundings, and apply the sign conventions for
        heat and work.
  \item Use the first law, $\Delta E = q + w$, including $w = -P\Delta V$.
  \item Explain what makes a quantity a state function.
\end{objectives}

\reading{Zumdahl \S6.1}{283--290}

\begin{warmup}[3]
  \item Convert \SI{25}{\celsius} to kelvin: \answerblank[2.6cm]{\SI{298}{\kelvin}}
  \item KMT: temperature measures what?
        \answerblank[4.6cm]{average kinetic energy}
  \item Moles in \SI{16.04}{\gram} \ce{CH4}: \answerblank[2.6cm]{\SI{1.000}{\mole}}
\end{warmup}

\segment{INSTRUCTION A}{25}{Energy, heat, and work}

\notesectionz{Defining the boundary}{6.1}
\practice{1}

Before any energy bookkeeping, you must decide what you are keeping books
\emph{on}:

\begin{itemize}[leftmargin=*,itemsep=0.2em]
  \item The \term{system} is the part of the universe you are studying ---
        usually the \answerblank[2.4cm]{reaction} itself.
  \item The \term{surroundings} are \answerblank[3.6cm]{everything else} ---
        the solvent, the container, the room.
\end{itemize}

Energy crosses that boundary in exactly two ways:

\renewcommand{\arraystretch}{1.4}
\noindent\begin{tabular}{@{}p{2.4cm}p{4.8cm}p{6.0cm}@{}}
\toprule
 & \textbf{Definition} & \textbf{Sign convention (system's view)}\\
\midrule
\term{Heat} ($q$) & transfer driven by a
  \answerblank[4.7cm]{temperature difference} &
  $q > 0$: heat flows \answerblank[1.6cm]{into} the system (endothermic)\newline
  $q < 0$: heat flows \answerblank[1.6cm]{out} (exothermic)\\
\term{Work} ($w$) & transfer by force acting through a distance &
  $w > 0$: work done \answerblank[2.2cm]{on} the system\newline
  $w < 0$: system does work on the surroundings\\
\bottomrule
\end{tabular}

\begin{apcnote}
Every sign is written from the \emph{system's} point of view --- treat the
system as an account and ask whether energy is being deposited or withdrawn.
Heat flowing out of a reaction is negative for the reaction even though the
room got warmer. Students who reason from the surroundings invert every
sign.
\end{apcnote}

\notesub{How heat transfer actually happens}

Zumdahl's picture: two gas samples at different temperatures separated by a
thin membrane. Temperature measures average kinetic energy, so the
fast-moving particles on the hot side transfer energy through collisions to
the slower particles on the cold side. Energy always flows from
\answerblank[1.6cm]{hot} to \answerblank[1.6cm]{cold}, until both reach the
same \answerblank[2.6cm]{intermediate} temperature.

\segment{GUIDED PRACTICE}{15}{Assign the signs}

\begin{enumerate}[leftmargin=*,itemsep=0.3em]
  \item A reaction warms the beaker it is in.
        \answerblank[5cm]{exothermic, $q < 0$}
  \item An ice pack cools your skin.
        \answerblank[6.4cm]{endothermic, $q > 0$ for the pack}
  \item A gas is compressed by a piston.
        \answerblank[5.9cm]{$w > 0$ (work done on the gas)}
  \item A gas expands and pushes the piston out.
        \answerblank[5cm]{$w < 0$}
\end{enumerate}

\segment{INSTRUCTION B}{20}{The first law}

\notesectionz{Conservation of energy, written down}{6.1}
\practice{5}

\[
  \Delta E = \answerblank[2cm]{$q + w$}
\]

$\Delta E$ is the change in \term{internal energy} --- the total kinetic and
potential energy of everything in the system. The first law says energy is
neither created nor destroyed; it only moves or changes form.

\notesub{Pressure--volume work}

When a gas expands against a constant external pressure:
\[
  w = \answerblank[2.4cm]{$-P\Delta V$}
\]

The minus sign is not decoration. If the gas \emph{expands}, $\Delta V$ is
positive, so $w$ is \answerblank[1.8cm]{negative} --- the system spent
energy pushing the surroundings back. If it is compressed, $\Delta V < 0$
and $w > 0$.

\begin{workedexample}[Worked example 1: combining heat and work]
A gas absorbs \SI{450}{\joule} of heat and expands, doing \SI{120}{\joule}
of work on its surroundings.
\[
  \Delta E = q + w = (+450) + (-120) = \SI{+330}{\joule}
\]
The system gained 450 J but spent 120 J pushing outward, so its internal
energy rose by only 330 J.
\end{workedexample}

\begin{workedexample}[Worked example 2: $P\Delta V$ with units]
A gas expands from \SI{2.0}{\liter} to \SI{5.0}{\liter} against a constant
\SI{1.5}{\atm}.
\[
  w = -P\Delta V = -(1.5)(5.0-2.0) = \SI{-4.5}{\liter\atm}
\]
Converting with \SI{1}{\liter\atm} $=$ \SI{101.3}{\joule}:
$w = \SI{-456}{\joule}$. Negative, as expected for expansion.
\end{workedexample}

\segment{APPLICATION}{20}{State functions}

\notesectionz{Path independence}{6.1}
\practice{6}

A \term{state function} depends only on the current state, not on how you
got there. Energy ($E$) and enthalpy ($H$) are state functions;
$q$ and $w$ individually are \answerblank[1.6cm]{not}.

\begin{apcnote}
Altitude is a state function; distance walked is not. Two hikers reaching
the same summit have identical altitude change but may have walked wildly
different distances. Likewise, two routes from the same reactants to the
same products give identical $\Delta E$ --- even though the heat and work
split differently along the way. Block 3 turns this fact into a calculation
tool.
\end{apcnote}

\begin{enumerate}[leftmargin=*,itemsep=0.4em]
  \item A system releases \SI{200}{\joule} of heat while
        \SI{75}{\joule} of work is done on it. Find $\Delta E$.
        \answerblank[3.2cm]{$-125$ J}
  \item Explain why $\Delta E$ for converting graphite to diamond is the
        same whether done in one step or ten.
        \answerlines[2]{Internal energy is a state function --- it depends
        only on the initial and final states, so any path connecting them
        gives the same $\Delta E$.}
\end{enumerate}

\begin{exitticket}
A reaction in a beaker makes the beaker feel cold. State the sign of $q$ for
the reaction and name the process. \answerlines[2]{$q > 0$ --- endothermic.
The reaction absorbed energy from the beaker and solution, lowering their
temperature.}
\end{exitticket}

% =====================================================================
\blocklesson{Enthalpy and Calorimetry}{Zumdahl \S6.2}

\begin{objectives}
  \item Relate enthalpy change to heat at constant pressure.
  \item Use $q = mc\Delta T$ and interpret calorimetry data.
\end{objectives}

\reading{Zumdahl \S6.2}{290--297}

\begin{warmup}[4]
  \item Sign of $q$ for an exothermic reaction:
        \answerblank[2cm]{negative}
  \item First law equation: \answerblank[2.4cm]{$\Delta E = q + w$}
  \item $w$ when a gas expands: \answerblank[2.4cm]{negative}
  \item Is enthalpy a state function? \answerblank[2cm]{yes}
\end{warmup}

\segment{INSTRUCTION A}{25}{Enthalpy}

\notesectionz{Why chemists use $H$ instead of $E$}{6.2}
\practice{5}

Enthalpy is defined as $H = E + PV$. Its value comes from one consequence:
at \emph{constant pressure} --- which is how nearly every reaction in an
open beaker runs ---
\[
  \Delta H = \answerblank[1.6cm]{$q_P$}
\]

So the heat you measure in an open container \emph{is} the enthalpy change.
That is the whole reason enthalpy exists as a bookkeeping quantity.

\renewcommand{\arraystretch}{1.3}
\noindent\begin{tabular}{@{}p{3.4cm}p{4.4cm}p{5.6cm}@{}}
\toprule
 & \textbf{$\Delta H$ sign} & \textbf{Surroundings}\\
\midrule
\term{Exothermic} & \answerblank[1.9cm]{negative} &
  get \answerblank[1.6cm]{warmer}\\
\term{Endothermic} & \answerblank[1.6cm]{positive} &
  get \answerblank[1.6cm]{cooler}\\
\bottomrule
\end{tabular}

\notesub{$\Delta H$ scales with amount}

A thermochemical equation's $\Delta H$ belongs to the coefficients as
written:
\[ \ce{CH4(g) + 2 O2(g) -> CO2(g) + 2 H2O(l)} \qquad \Delta H = \SI{-890}{\kilo\joule} \]

\begin{workedexample}[Worked example 1: scaling]
How much heat is released when \SI{5.8}{\gram} of \ce{CH4} burns at constant
pressure?
\begin{align*}
  n &= \frac{5.8}{16.04} = \SI{0.36}{\mole}\\
  \Delta H &= 0.36 \times (-890) = \SI{-320}{\kilo\joule}
\end{align*}
\emph{Reality check:} less than one mole burned, so less than 890 kJ should
be released. It is.
\end{workedexample}

\segment{GUIDED PRACTICE}{15}{Enthalpy scaling}

Using $\Delta H = \SI{-890}{\kilo\joule}$ per mole of \ce{CH4}:
\begin{enumerate}[leftmargin=*,itemsep=0.3em]
  \item Heat from \SI{2.00}{\mole} \ce{CH4}:
        \answerblank[3cm]{\SI{-1780}{\kilo\joule}}
  \item Heat from \SI{32.08}{\gram} \ce{CH4}:
        \answerblank[3cm]{\SI{-1780}{\kilo\joule}}
  \item $\Delta H$ for the \emph{reverse} reaction:
        \answerblank[3cm]{\SI{+890}{\kilo\joule}}
\end{enumerate}

\segment{INSTRUCTION B}{20}{Calorimetry}

\notesectionz{Measuring heat by measuring temperature}{6.2}
\practice{4}

A \term{calorimeter} determines heat by watching a temperature change. The
central equation:
\[
  q = \answerblank[2.2cm]{$mc\Delta T$}
  \qquad \Delta T = T_{\text{final}} - T_{\text{initial}}
\]

The \term{specific heat capacity} $c$ is the energy needed to raise
\SI{1}{\gram} by \SI{1}{\celsius}. For water,
$c = \answerblank[3.2cm]{\SI{4.184}{\joule\per\gram\per\celsius}}$ --- unusually
large, which is why water moderates temperature so effectively.

\begin{apctrap}
The heat absorbed by the \emph{solution} and the heat released by the
\emph{reaction} are equal in magnitude and opposite in sign:
$q_{\text{rxn}} = -q_{\text{solution}}$. Students routinely report an
exothermic reaction with a positive $\Delta H$ because they forgot to flip
the sign after computing $mc\Delta T$ for the water.
\end{apctrap}

\begin{workedexample}[Worked example 2: neutralization in a coffee cup]
\SI{50.0}{\milli\liter} of \SI{1.00}{\Molar} \ce{HCl} and
\SI{50.0}{\milli\liter} of \SI{1.00}{\Molar} \ce{NaOH}, both at
\SI{25.0}{\celsius}, are mixed. The temperature rises to
\SI{31.9}{\celsius}. Assume the solution has the density and specific heat
of water.
\begin{align*}
  m &= \SI{100.0}{\gram}, \qquad \Delta T = \SI{6.9}{\celsius}\\
  q_{\text{soln}} &= (100.0)(4.184)(6.9) = \SI{2887}{\joule} = \SI{2.89}{\kilo\joule}\\
  q_{\text{rxn}} &= \SI{-2.89}{\kilo\joule}\\
  n &= (0.0500)(1.00) = \SI{0.0500}{\mole}\\
  \Delta H &= \frac{-2.89}{0.0500} = \SI{-57.8}{\kilo\joule\per\mole}
\end{align*}
The accepted value is about \SI{-57.3}{\kilo\joule\per\mole} --- close, with
the difference explained by heat lost to the cup and thermometer.
\end{workedexample}

\segment{APPLICATION}{20}{Determining a specific heat}

A \SI{55.0}{\gram} piece of metal is heated to \SI{99.0}{\celsius} and
dropped into \SI{100.0}{\gram} of water at \SI{21.0}{\celsius}. The final
temperature is \SI{24.8}{\celsius}.

\begin{enumerate}[leftmargin=*,itemsep=0.4em]
  \item Heat gained by the water: \workspace[1.8cm]
        \keyonly{\begin{apcanswer}$q = (100.0)(4.184)(3.8) =
        \SI{1590}{\joule}$\end{apcanswer}}
  \item Specific heat of the metal: \workspace[2.2cm]
        \keyonly{\begin{apcanswer}$q_{\text{metal}} = \SI{-1590}{\joule}$;
        $c = 1590/[(55.0)(74.2)] =
        \SI{0.390}{\joule\per\gram\per\celsius}$\end{apcanswer}}
  \item Suggest the metal's identity.
        \answerblank[4cm]{copper ($c \approx 0.385$)}
  \item Why does the water's temperature rise far less than the metal's
        falls? \answerlines[2]{Water's specific heat is about eleven times
        larger, and there is nearly twice the mass --- so absorbing the same
        energy produces a much smaller temperature change.}
\end{enumerate}

\begin{exitticket}
A student computes $q_{\text{solution}} = \SI{+3.2}{\kilo\joule}$ for a
reaction and reports $\Delta H = \SI{+3.2}{\kilo\joule}$. What did they get
wrong? \answerlines[2]{The solution \emph{gained} heat, so the reaction
\emph{released} it: $q_{\text{rxn}} = \SI{-3.2}{\kilo\joule}$. The reaction
is exothermic and $\Delta H$ is negative.}
\end{exitticket}

% =====================================================================
\blocklesson{Hess's Law}{Zumdahl \S6.3}

\begin{objectives}
  \item State Hess's law and explain why it follows from $H$ being a state
        function.
  \item Combine thermochemical equations to obtain an unknown $\Delta H$.
\end{objectives}

\reading{Zumdahl \S6.3}{297--301}

\begin{warmup}[3]
  \item $\Delta H$ equals $q$ under what condition?
        \answerblank[3.4cm]{constant pressure}
  \item $q$ for \SI{25.0}{\gram} water warmed \SI{10.0}{\celsius}:
        \answerblank[3cm]{\SI{1046}{\joule}}
  \item Sign of $\Delta H$ for an endothermic process:
        \answerblank[2.2cm]{positive}
\end{warmup}

\segment{INSTRUCTION A}{25}{The law and why it must be true}

\notesectionz{Hess's law}{6.3}
\practice{5}

\begin{apcnote}
\textbf{Hess's law:} if a reaction is carried out in a series of steps,
$\Delta H$ for the overall reaction equals the \answerblank[2cm]{sum} of the
$\Delta H$ values for the individual steps.

This is not a separate discovery --- it is a direct consequence of enthalpy
being a \answerblank[3.0cm]{state function}. If $\Delta H$ depended on the
route, you could build a machine that created energy by going one way and
returning another.
\end{apcnote}

\begin{workedexample}[Worked example 1: Zumdahl's nitrogen oxides]
The direct reaction:
\[ \ce{N2(g) + 2 O2(g) -> 2 NO2(g)} \qquad \Delta H_1 = \SI{68}{\kilo\joule} \]
The same change in two steps:
\begin{align*}
  \ce{N2(g) + O2(g) &-> 2 NO(g)}   &\Delta H_2 &= \SI{180}{\kilo\joule}\\
  \ce{2 NO(g) + O2(g) &-> 2 NO2(g)} &\Delta H_3 &= \SI{-112}{\kilo\joule}
\end{align*}
Adding the steps cancels the \ce{2 NO} and reproduces the overall equation,
and indeed $180 + (-112) = \SI{68}{\kilo\joule} = \Delta H_1$.
\end{workedexample}

\notesub{The two manipulation rules}

\begin{enumerate}[leftmargin=*,itemsep=0.2em]
  \item Reverse a reaction $\Rightarrow$ \answerblank[3.3cm]{reverse the sign}
        of $\Delta H$.
  \item Multiply the coefficients by a factor $\Rightarrow$
        \answerblank[3.4cm]{multiply $\Delta H$} by the same factor.
\end{enumerate}

Rule 1 follows because $\Delta H$'s sign records the direction of heat flow;
run the change backwards and the flow reverses. Rule 2 follows because
enthalpy is \term{extensive} --- twice the reaction releases twice the heat.

\segment{GUIDED PRACTICE}{15}{Manipulate}

Given \ce{2 H2(g) + O2(g) -> 2 H2O(l)}, $\Delta H = \SI{-572}{\kilo\joule}$:
\begin{enumerate}[leftmargin=*,itemsep=0.3em]
  \item $\Delta H$ for \ce{H2(g) + 1/2 O2(g) -> H2O(l)}:
        \answerblank[3cm]{\SI{-286}{\kilo\joule}}
  \item $\Delta H$ for \ce{2 H2O(l) -> 2 H2(g) + O2(g)}:
        \answerblank[3cm]{\SI{+572}{\kilo\joule}}
  \item $\Delta H$ for \ce{4 H2(g) + 2 O2(g) -> 4 H2O(l)}:
        \answerblank[3cm]{\SI{-1144}{\kilo\joule}}
\end{enumerate}

\segment{INSTRUCTION B}{20}{The working method}

\notesectionz{A reliable procedure}{6.3}
\practice{5}

\begin{enumerate}[leftmargin=*,itemsep=0.2em]
  \item Write the \answerblank[3.3cm]{target equation} first.
  \item Find a substance that appears in only \emph{one} given equation and
        place it correctly --- reverse or scale that equation as needed.
  \item Repeat for the next such substance.
  \item Add everything; confirm the intermediates
        \answerblank[1.6cm]{cancel}.
  \item Sum the adjusted $\Delta H$ values.
\end{enumerate}

\begin{workedexample}[Worked example 2: a two-step target]
Target: \ce{C(s) + 1/2 O2(g) -> CO(g)}, given
\begin{align*}
  \text{(i)}\ \ce{C(s) + O2(g) &-> CO2(g)} &\Delta H &= \SI{-393.5}{\kilo\joule}\\
  \text{(ii)}\ \ce{2 CO(g) + O2(g) &-> 2 CO2(g)} &\Delta H &= \SI{-566.0}{\kilo\joule}
\end{align*}
CO appears only in (ii), but on the wrong side and with the wrong
coefficient. Reverse (ii) and halve it:
\[ \ce{CO2(g) -> CO(g) + 1/2 O2(g)} \qquad \Delta H = \SI{+283.0}{\kilo\joule} \]
Add to (i): the \ce{CO2} cancels, giving the target with
\[ \Delta H = -393.5 + 283.0 = \SI{-110.5}{\kilo\joule} \]
\end{workedexample}

\segment{APPLICATION}{20}{Build it yourself}

Determine $\Delta H$ for \ce{2 C(s) + H2(g) -> C2H2(g)} given:
\begin{align*}
  \text{(i)}\ \ce{C2H2(g) + 5/2 O2(g) &-> 2 CO2(g) + H2O(l)} &\Delta H &= \SI{-1300.}{\kilo\joule}\\
  \text{(ii)}\ \ce{C(s) + O2(g) &-> CO2(g)} &\Delta H &= \SI{-393.5}{\kilo\joule}\\
  \text{(iii)}\ \ce{H2(g) + 1/2 O2(g) &-> H2O(l)} &\Delta H &= \SI{-286}{\kilo\joule}
\end{align*}

\workspace[4.4cm]
\keyonly{\begin{apcanswer}Reverse (i): $+1300.$; use (ii) $\times 2$:
$-787.0$; use (iii) as written: $-286$.
$\Delta H = 1300. - 787.0 - 286 = \SI{+227}{\kilo\joule}$ --- acetylene is
endothermic to form, which is why it stores so much energy.\end{apcanswer}}

\begin{exitticket}
Why does reversing a reaction reverse the sign of $\Delta H$?
\answerlines[2]{The sign records the direction of heat flow at constant
pressure. Running the change in reverse reverses that flow --- energy
released one way must be absorbed the other.}
\end{exitticket}

% =====================================================================
\blocklesson{Standard Enthalpies of Formation}{Zumdahl \S6.4}

\begin{objectives}
  \item Define $\Delta H^\circ_f$ and standard state.
  \item Compute $\Delta H^\circ_{\text{rxn}}$ from tabulated formation
        enthalpies.
\end{objectives}

\reading{Zumdahl \S6.4}{301--308}

\begin{warmup}[4]
  \item Hess's law depends on $H$ being a what?
        \answerblank[3cm]{state function}
  \item Reverse a reaction: $\Delta H$ does what?
        \answerblank[3cm]{changes sign}
  \item Halve the coefficients: $\Delta H$ does what?
        \answerblank[2.4cm]{halves}
  \item $q$ for \SI{50.0}{\gram} water cooled \SI{4.00}{\celsius}:
        \answerblank[3cm]{\SI{-837}{\joule}}
\end{warmup}

\segment{INSTRUCTION A}{25}{A universal reference point}

\notesectionz{Standard enthalpy of formation}{6.4}
\practice{5}

The \term{standard enthalpy of formation} $\Delta H^\circ_f$ is the enthalpy
change when \answerblank[1.4cm]{1 mole} of a compound forms from its
\answerblank[2cm]{elements}, with everything in its
\term{standard state}.

Standard state means the substance's most stable form at
\SI{1}{\atm} and (conventionally) \SI{25}{\celsius} --- so
\ce{O2(g)}, \ce{Br2(l)}, \ce{C(graphite)}.

\begin{apcnote}
It follows directly from the definition that
$\Delta H^\circ_f$ of an \emph{element in its standard state} is
\answerblank[1.2cm]{0} --- forming an element from itself is no change at
all. That is why \ce{O2(g)} contributes nothing to a calculation while
\ce{O3(g)} (not the standard form) contributes $+143$ kJ/mol.
\end{apcnote}

\notesub{Why the tables exist}

Some enthalpy changes cannot be measured directly. Zumdahl's example:
converting graphite to diamond is far too slow for a calorimeter. But if
every compound's formation enthalpy is tabulated once, any reaction's
$\Delta H$ can be assembled from them --- Hess's law industrialized.

\segment{GUIDED PRACTICE}{15}{Writing formation equations}

Write the formation equation (1 mol of product, elements in standard states):
\begin{enumerate}[leftmargin=*,itemsep=0.35em]
  \item \ce{NH3(g)}:
        \answerblank[6cm]{\ce{1/2 N2(g) + 3/2 H2(g) -> NH3(g)}}
  \item \ce{CO2(g)}:
        \answerblank[6.7cm]{\ce{C(graphite) + O2(g) -> CO2(g)}}
  \item $\Delta H^\circ_f$ of \ce{N2(g)}: \answerblank[2cm]{0}
\end{enumerate}

\segment{INSTRUCTION B}{20}{The master equation}

\notesectionz{Computing $\Delta H^\circ_{\text{rxn}}$}{6.4}
\practice{5}

\[
  \Delta H^\circ_{\text{rxn}} =
  \answerblank[8.0cm]{$\sum n\,\Delta H^\circ_f(\text{products}) - \sum n\,\Delta H^\circ_f(\text{reactants})$}
\]

Note the two easy losses: multiply each value by its
\answerblank[2cm]{coefficient}, and subtract in the right order ---
\answerblank[5.2cm]{products minus reactants}.

\begin{workedexample}[Worked example 1: methane combustion]
\[ \ce{CH4(g) + 2 O2(g) -> CO2(g) + 2 H2O(l)} \]
Using $\Delta H^\circ_f$: \ce{CH4} $-74.8$, \ce{O2} $0$, \ce{CO2} $-393.5$,
\ce{H2O(l)} $-285.8$ kJ/mol.
\begin{align*}
  \Delta H^\circ &= [(-393.5) + 2(-285.8)] - [(-74.8) + 2(0)]\\
  &= (-965.1) - (-74.8) = \SI{-890.3}{\kilo\joule}
\end{align*}
This matches the experimental $-890$ kJ used in Block 2 --- a good check
that the method and the table agree.
\end{workedexample}

\begin{apctrap}
Watch the physical state. $\Delta H^\circ_f$ for \ce{H2O(l)} is $-285.8$ but
for \ce{H2O(g)} it is $-241.8$ --- a 44 kJ/mol difference, exactly the heat
of vaporization. Using the wrong state is worth 88 kJ of error in the
example above.
\end{apctrap}

\segment{APPLICATION}{20}{Full calculations}

\begin{enumerate}[leftmargin=*,itemsep=0.4em]
  \item Compute $\Delta H^\circ$ for
        \ce{2 NH3(g) -> N2(g) + 3 H2(g)}, given
        $\Delta H^\circ_f(\ce{NH3}) = \SI{-46.1}{\kilo\joule\per\mole}$.
        \workspace[2.2cm]
        \keyonly{\begin{apcanswer}$\Delta H^\circ = [0 + 0] - [2(-46.1)]
        = \SI{+92.2}{\kilo\joule}$ --- decomposition is
        endothermic.\end{apcanswer}}
  \item Compute $\Delta H^\circ$ for the thermite reaction
        \ce{2 Al(s) + Fe2O3(s) -> Al2O3(s) + 2 Fe(s)}, given
        $\Delta H^\circ_f$: \ce{Fe2O3} $-826$, \ce{Al2O3} $-1676$
        kJ/mol. \workspace[2.4cm]
        \keyonly{\begin{apcanswer}$\Delta H^\circ = [(-1676) + 0] -
        [0 + (-826)] = \SI{-850}{\kilo\joule}$ --- strongly exothermic,
        which is why thermite melts iron.\end{apcanswer}}
  \item In problem 2, why do \ce{Al(s)} and \ce{Fe(s)} contribute nothing?
        \answerlines[2]{Both are elements in their standard states, so their
        $\Delta H^\circ_f$ values are zero by definition.}
\end{enumerate}

\begin{exitticket}
A reaction has $\Delta H^\circ = \SI{-1250}{\kilo\joule}$. State whether it
is exothermic or endothermic, and what happens to the temperature of the
surroundings. \answerlines[2]{Exothermic --- the negative sign means heat
flows out of the system, so the surroundings warm up.}
\end{exitticket}

\end{document}
